% Shared preamble of the HSEMHS mode documents (docs/physics/*.tex, docs/technical/*.tex).
\documentclass[10pt]{extarticle}
\usepackage[a4paper,margin=1.7cm]{geometry}
\usepackage[T1]{fontenc}\usepackage{lmodern}\usepackage{tabularx}
\usepackage{amsmath,amssymb,bm,enumitem,titlesec,xcolor,booktabs,array,tikz,graphicx,longtable}
\usepackage[hidelinks]{hyperref}
\titleformat{\section}{\large\bfseries\color{blue!50!black}}{\thesection.}{0.5em}{}
\titlespacing*{\section}{0pt}{9pt}{3pt}
\titleformat{\subsection}{\normalsize\bfseries}{\thesubsection}{0.5em}{}
\titlespacing*{\subsection}{0pt}{6pt}{2pt}
\setlist{nosep,leftmargin=1.3em}
\setlength{\parindent}{0pt}\setlength{\parskip}{4pt}
\renewcommand{\arraystretch}{1.15}
\newcommand{\B}{\mathbf B}\newcommand{\J}{\mathbf J}\newcommand{\F}{\mathbf F}\newcommand{\R}{\mathbf R}
\newcommand{\zh}{\hat{\mathbf z}}\newcommand{\Bo}{\mathbf B_{\rm obs}}\newcommand{\Lv}{\mathbf L}
\newcommand{\code}[1]{\texttt{#1}}
\newcolumntype{L}[1]{>{\raggedright\arraybackslash}p{#1}}
\newcolumntype{Y}{>{\raggedright\arraybackslash}X}
\newcommand{\doctitle}[2]{{\Large\bfseries #1}\\[2pt]{\small #2}\par\vspace{2pt}}
\pagestyle{plain}
\AtBeginDocument{\sloppy}

\begin{document}
\doctitle{column-hse-z: hydrostatic columns on a geometric-height grid}{Physics and method of \code{run.mode = "column-hse-z"} in HSEMHS. A self-contained description.}

\begin{abstract}
\noindent Given the temperature of a three-dimensional piece of the solar atmosphere on a grid of geometric heights, this
mode computes the gas pressure, the density, the electron density, the 500\,nm extinction and the optical depth at every
grid point. It does so by treating every vertical column of the grid as an isolated, static column of gas that carries its
own weight (hydrostatic equilibrium), with the same prescribed pressure at one level of every column. The equation of
state includes ionisation and molecules. The equation solved, its boundary condition, the numerical integration, the
accuracy control, the computation of the optical depth and the assumptions are derived and explained step by step, as
well as what the result does not satisfy: horizontal force balance.
\end{abstract}

\section{Introduction}
A radiative magnetohydrodynamic simulation or a model atmosphere describes the gas by its temperature $T$, its density
$\rho$, its pressure $P$ and its velocity. If only the temperature is kept, the simplest physically consistent way to
recover the pressure and the density is to assume that the gas is \emph{at rest} and that each vertical column supports its
own weight. This mode implements exactly that. Its output serves two purposes: it is a stand-alone hydrostatic
atmosphere with an optical-depth scale, and it is the starting point of the magnetohydrostatic modes, which add the
horizontal force balance and the magnetic field.

\textbf{Input:} the temperature $T(x,y,z)$ on a rectangular grid, the coordinates $x$, $y$, $z$, the gravity $g$ and one
reference pressure $P_{\rm ref}$. \textbf{Output:} $P$, $\rho$, $n_e$, $P_e$, $\bar\mu$, $n_{\rm H}$, the column mass $m$,
$\chi_{500}$ and $\tau_{500}$ on the same grid, the height of $\tau_{500}=1$ in every column, and a diagnostic of the
horizontal and vertical force imbalance.

\section{Notation}
\begin{tabularx}{\linewidth}{>{$}l<{$} Y l l}
\toprule
\text{Symbol} & Meaning & Unit & Role \\ \midrule
x,\ y & horizontal coordinates; grid $x_i$, $i=1\dots N_x$, $y_j$, $j=1\dots N_y$, uniform spacing $\Delta x$, $\Delta y$ & m & given \\
z & height, growing upward; grid $z_k$, $k=1\dots N_z$, stored top first ($z_1>z_2>\dots>z_{N_z}$) & m & given \\
T & temperature & K & given, fixed \\
X & $\ln(T/1\,{\rm K})$ & -- & given \\
P & gas pressure & Pa & unknown \\
u & $\ln(P/1\,{\rm Pa})$, the variable that is integrated & -- & unknown \\
\rho & mass density, $\rho(T,P)$ from the equation of state & kg\,m$^{-3}$ & derived \\
q & $\rho/P$ & s$^2$\,m$^{-2}$ & derived \\
n_e,\ P_e & electron number density; electron pressure $n_ek_BT$ & m$^{-3}$, Pa & derived \\
\bar\mu & mean molecular weight, $\rho k_BT/(P m_{\rm u})$ & -- & derived \\
n_{\rm H} & number density of hydrogen nuclei, $\rho$ divided by the mass per H nucleus & m$^{-3}$ & derived \\
H & pressure scale height $P/(\rho g)$ & m & derived \\
m & column mass: mass above a point per unit horizontal area & kg\,m$^{-2}$ & derived \\
g & gravitational acceleration (274\,m\,s$^{-2}$ at the solar surface) & m\,s$^{-2}$ & parameter \\
P_{\rm ref} & pressure prescribed at the anchor level of every column & Pa & parameter \\
z_{\rm ref} & anchor height (only for the anchor option \code{height}) & m & parameter \\
\chi_{500} & extinction coefficient at 500\,nm (absorption plus scattering) & m$^{-1}$ & derived \\
\tau_{500} & optical depth at 500\,nm along the vertical, counted downward from the top & -- & derived \\
\tau_{\rm top} & optical depth of the gas above the top of the grid & -- & derived \\
H_\chi & scale height of $\chi_{500}$ at the top of the grid & m & derived \\
\mu & cosine of the viewing angle (default 1, vertical) & -- & parameter \\
k_B,\ m_{\rm u} & Boltzmann constant $1.381\times10^{-23}$\,J\,K$^{-1}$; atomic mass unit $1.661\times10^{-27}$\,kg & & constants \\
\F & net force per unit volume & N\,m$^{-3}$ & diagnostic \\
\bottomrule
\end{tabularx}

\section{The force balance of a static gas}\label{sec:force}
\subsection{All three components}
A small volume of gas is at rest when the forces on it cancel. Without a magnetic field, two forces act per unit volume:
the pressure force $-\nabla P$ (gas is pushed from high to low pressure) and gravity $-\rho g\,\zh$. The static momentum
equation is therefore
\begin{equation}
\F\equiv-\nabla P-\rho\,g\,\zh=0,\qquad\text{i.e.}\qquad
-\frac{\partial P}{\partial x}=0,\quad -\frac{\partial P}{\partial y}=0,\quad -\frac{\partial P}{\partial z}-\rho g=0 .
\label{eq:F}
\end{equation}
This mode imposes \emph{only the vertical component}, separately in every column. The two horizontal components are not
imposed; Sect.~\ref{sec:3d} shows why they cannot be satisfied in general, and how large they remain is reported as a
diagnostic.

\subsection{Hydrostatic equilibrium of a column}
Take a thin horizontal slab of area $A$ between $z$ and $z+dz$. The gas below pushes it up with the force $P(z)A$, the gas
above pushes it down with $P(z+dz)A$, and gravity pulls it down with its weight $\rho\,A\,dz\,g$. Equilibrium means
$P(z)A-P(z+dz)A-\rho gA\,dz=0$. Dividing by $A\,dz$ and letting $dz\to0$ gives the \textbf{hydrostatic equation}
\begin{equation}
\boxed{\;\frac{dP}{dz}=-\rho\,g\;}\label{eq:hse}
\end{equation}
which is the $z$ component of Eq.~\eqref{eq:F}: the pressure falls upward, fastest where the gas is densest.

% Shared section of the physics documents: equation of state and 500 nm extinction (\input by each mode document).
\section{The equation of state and the opacity}\label{sec:eos}
Every force balance in this document involves the density $\rho$, and every optical depth involves the extinction
coefficient $\chi_{500}$. Neither is an unknown of its own: at given temperature and pressure both are fixed by the
microphysics of the gas. This section explains how.

\subsection{The ideal gas}
In the solar photosphere and chromosphere the particles interact only weakly, so the gas obeys the ideal-gas law. If $n$
is the number of free particles (atoms, ions, molecules and electrons) per cubic metre,
\begin{equation}
P=n\,k_BT .\label{eq:ideal}
\end{equation}
If the average mass of a free particle is $\bar\mu\,m_{\rm u}$ (with $m_{\rm u}$ the atomic mass unit and $\bar\mu$ the
\textbf{mean molecular weight}), the mass density is $\rho=n\,\bar\mu\,m_{\rm u}$, and therefore
\begin{equation}
\rho=\frac{\bar\mu\,m_{\rm u}}{k_B}\,\frac{P}{T}.\label{eq:rho_ideal}
\end{equation}
At fixed $\bar\mu$ and $T$ the density is proportional to the pressure.

\subsection{\texorpdfstring{Why $\bar\mu$ depends on $T$ and $P$}{Why the mean molecular weight depends on T and P}: ionisation and molecules}
For neutral solar gas (about 91\,\% hydrogen and 9\,\% helium atoms by number, plus a small fraction of heavier elements)
$\bar\mu\approx1.3$. When a hydrogen atom is ionised it becomes two free particles, a proton and an electron, of almost
the same total mass; in fully ionised gas $\bar\mu\approx0.6$. In the coolest layers hydrogen partly forms H$_2$ molecules,
which raises $\bar\mu$. The degree of ionisation of an element in stage $i$ is given by the \textbf{Saha equation}
\begin{equation}
\frac{n_{i+1}\,n_e}{n_i}=\frac{2U_{i+1}(T)}{U_i(T)}\left(\frac{2\pi m_ek_BT}{h^2}\right)^{3/2}e^{-E_i/k_BT},\label{eq:saha}
\end{equation}
with $E_i$ the ionisation energy (13.6\,eV for hydrogen), $U_i$ the partition functions, $m_e$ the electron mass and $h$
Planck's constant. Two properties matter for us. (i) Because of the exponential, ionisation switches on over a narrow
range of temperature. (ii) The electron density $n_e$ appears on the left: at higher density there are more electrons to
recombine with, and a higher temperature is needed for the same ionisation. Hence $\bar\mu=\bar\mu(T,P)$, and
Eq.~\eqref{eq:rho_ideal} is no longer a simple proportionality. The Saha equations of all elements are coupled through the
common $n_e$ and through charge conservation ($n_e$ equals the total charge of the ions), and the molecular equilibria
(H$_2$, H$_2^+$, H$^-$) are solved together with them. This assumes \textbf{local thermodynamic equilibrium} (LTE): the
level populations and ionisation stages are those of a gas in equilibrium at the local $T$.

\subsection{\texorpdfstring{The 500\,nm extinction}{The 500 nm extinction}}
Light of wavelength 500\,nm travelling a distance $ds$ loses the fraction $\chi_{500}\,ds$ of its intensity, by absorption
and by scattering. In the solar photosphere $\chi_{500}$ is dominated by the negative hydrogen ion H$^-$ (a hydrogen atom
with a second, weakly bound electron, binding energy 0.75\,eV), whose abundance is proportional to the electron density.
The electrons come mostly from easily ionised metals (Na, Mg, Si, Ca, Fe, \dots), so $\chi_{500}$ grows steeply with
temperature and with density. The code adds the continuum contributions of H (bound--free and free--free), H$^-$,
H$_2^+$, He, He$^-$, the metals C, N, Mg, Si, Ca and Fe, Rayleigh scattering on H, He and H$_2$, and Thomson scattering
on free electrons, each counted once. Spectral lines are not included: $\chi_{500}$ is a continuum extinction.

\subsection{The table and its interpolation}
Solving the coupled Saha, molecular and charge equations at every grid point at every step would be far too slow. They
are therefore solved once, with Wittmann's equation-of-state code for solar composition (with a dedicated hot branch above
a switch temperature $T_{\rm sw}$, where the original solver fails in the hot, low-pressure corner), on a grid of
$X=\ln T$ and $Y=\ln P$, and stored in a table. For every node the table holds the three functions
\[
\ln\rho(X,Y),\qquad \ln n_e(X,Y),\qquad \ln\chi_{500}(X,Y),
\]
together with their derivatives $\partial_X$, $\partial_Y$ and $\partial_X\partial_Y$. Between the nodes each function is
evaluated by \textbf{bicubic Hermite interpolation}: a cubic polynomial in each direction that matches the stored values
and derivatives at the four corners of the cell. The result is continuous with continuous first derivatives, which is
important for the Newton-type solvers of the magnetohydrostatic modes, and it provides the derivative
$\partial\ln\rho/\partial\ln P|_T$ that they need. The table used by all modes is \code{cache/eos\_table\_fullz.h5}.
Its range of $T$ and $P$ (its \emph{envelope}) covers the input atmospheres; a point outside the envelope is reported as
an equation-of-state failure, never extrapolated.

Two derived quantities appear repeatedly:
\begin{equation}
q\equiv\frac{\rho}{P}\quad[{\rm s^2\,m^{-2}}],\qquad H\equiv\frac{P}{\rho\,g}=\frac{1}{g\,q}=\frac{k_BT}{\bar\mu\,m_{\rm u}\,g}\quad[{\rm m}].
\label{eq:H}
\end{equation}
$H$ is the \textbf{pressure scale height}. For $T=6000$\,K and $\bar\mu=1.3$, $H\approx140$\,km.


\section{The equation solved}\label{sec:eq}
\subsection{Logarithmic form}
Divide Eq.~\eqref{eq:hse} by $P$ and use $u=\ln P$ (so $du=dP/P$) and $\rho=q\,P$:
\begin{equation}
\boxed{\;\frac{du}{dz}=-g\,q\big(T(z),u\big)=-g\,\exp\!\big(\ln\rho(X,u)-u\big)=-\frac{1}{H(T,P)}\;}\label{eq:dudz}
\end{equation}
In an isothermal gas of constant $\bar\mu$, $H$ is a constant and the solution is $P=P_0\,e^{-(z-z_0)/H}$: the pressure falls
by a factor $e$ every scale height. Integrating in $u$ instead of $P$ has two advantages: a column spans eight or more
orders of magnitude in $P$, but $u$ changes smoothly and by moderate amounts per grid interval, which is what a numerical
integrator handles best; and $P=e^u$ is positive whatever happens.

\subsection{The boundary condition: one pressure per column}\label{sec:anchor}
Equation~\eqref{eq:dudz} is a first-order ordinary differential equation for $u(z)$ in each column. It fixes how the
pressure changes with height, not its value: if $u(z)$ is a solution, a different starting value gives a different, equally
valid solution. One value per column is needed, and the mode prescribes the \textbf{same pressure $P_{\rm ref}$ at the
same level in every column}:
\begin{equation}
u(x_i,y_j,z_a)=\ln P_{\rm ref}\qquad\text{for all columns }(i,j),\label{eq:anchor}
\end{equation}
where the anchor level $z_a$ is chosen by \code{column\_hse.anchor}:
\begin{itemize}
\item \code{bottom} (default): $z_a=z_{N_z}$, the lowest level; the integration runs upward.
\item \code{top}: $z_a=z_1$; the integration runs downward ($P_{\rm ref}$ defaults to 0.1\,Pa if not set).
\item \code{height}: $z_a=z_{\rm ref}$ anywhere inside the grid; the integration runs upward and downward from it.
\end{itemize}
$P_{\rm ref}$ is either a number in the configuration or, with \code{P\_ref = "from\_input"}, the attribute
\code{P\_ref\_Pa} of the input file. For the production input this attribute is the median of the simulation's own
pressure over the bottom plane.

\textbf{Why the bottom?} The anchor condition~\eqref{eq:anchor} declares the anchor plane to be a surface of constant
pressure. Deep in the atmosphere, below the visible surface, the gas is dense, the relative temperature differences
between columns are smaller and the true pressure varies much less across a horizontal plane than higher up, so this is
the level at which the assumption is least wrong. All differences between the columns above it then come from their
temperature profiles. With a top anchor the photospheric pressure of each column would instead depend on the hot,
strongly structured chromosphere and transition region above it.

\textbf{Existence and uniqueness.} The right-hand side of Eq.~\eqref{eq:dudz} is a smooth function of $z$ and $u$ (the
table interpolation has continuous first derivatives), so for every column the equation with its anchor value has
exactly one solution, as long as $(T,P)$ stays inside the table.

\subsection{Statement of the problem}
For every column $(i,j)$, $i=1\dots N_x$, $j=1\dots N_y$, independently: \emph{find $u(z)$ on $z_{N_z}\le z\le z_1$ such
that $du/dz=-g\,q(T,u)$ and $u(z_a)=\ln P_{\rm ref}$; then compute $\rho$, $n_e$, $\chi_{500}$ from the table, the column
mass and the optical depth.} There is no coupling between columns: the problem consists of $N_xN_y$ independent
one-dimensional initial-value problems.

\section{Assumptions}\label{sec:assume}
\begin{enumerate}[label=\textbf{A\arabic*.},leftmargin=2.2em]
\item \textbf{Static gas.} All velocities are neglected; the momentum flux and the pressure of turbulent or wave motions are
not included.
\item \textbf{No magnetic force.} The field, if present in the input, is not used.
\item \textbf{Vertical balance only, column by column.} Only the $z$ component of $\F=0$ is imposed; the horizontal
components are left unbalanced (Sect.~\ref{sec:3d}).
\item \textbf{Equal pressure on the anchor plane.} Every column has $P=P_{\rm ref}$ at $z_a$ (Sect.~\ref{sec:anchor}). This is
the only external pressure information.
\item \textbf{Constant, vertical gravity.} $g$ is the same everywhere; the curvature of the Sun and the decrease of $g$ with
height (0.4\,\% over 3\,Mm) are neglected (plane-parallel geometry).
\item \textbf{LTE ideal gas of solar composition}, with ionisation and molecules from the Saha and molecular equilibria
(Sect.~\ref{sec:eos}).
\item \textbf{Temperature between grid levels:} $\ln T$ varies linearly with $z$ inside each grid interval.
\item \textbf{Optical depth:} continuum extinction at 500\,nm along the vertical (or along a fixed angle with $\mu\ne1$),
without lines; the gas above the grid is represented by an exponential extrapolation of $\chi_{500}$ (Sect.~\ref{sec:tau}).
\end{enumerate}

\section{Numerical solution}\label{sec:num}
\subsection{Temperature inside an interval}
The temperature is known only on the grid levels. Inside the interval $[z_{k+1},z_k]$ the integrator needs it at
arbitrary heights; by assumption A7,
\begin{equation}
X(z)=X_{k+1}+(X_k-X_{k+1})\,\frac{z-z_{k+1}}{z_k-z_{k+1}},\qquad X=\ln T .
\end{equation}
Interpolating $\ln T$ rather than $T$ keeps the temperature positive and represents an exponential rise (such as the
transition region) better.

\subsection{The classical Runge--Kutta method}
Write Eq.~\eqref{eq:dudz} as $du/dz=f(z,u)$ with $f=-g\,q(X(z),u)$. A step of length $s$ from $(z,u)$ evaluates the slope
four times,
\begin{equation}
\begin{aligned}
k_1&=f(z,\,u), & k_2&=f\!\left(z+\tfrac{s}{2},\,u+\tfrac{s}{2}k_1\right),\\
k_3&=f\!\left(z+\tfrac{s}{2},\,u+\tfrac{s}{2}k_2\right), & k_4&=f(z+s,\,u+s\,k_3),
\end{aligned}
\qquad
u(z+s)\approx u+\frac{s}{6}\,(k_1+2k_2+2k_3+k_4).\label{eq:rk4}
\end{equation}
The slopes at the beginning, twice in the middle and at the end are averaged with weights $1,2,2,1$ (Simpson's rule). The
error per step is proportional to $s^5$ and the total error to $s^4$: halving the step reduces the error by about a factor 16.
For an upward integration $s>0$, for a downward one $s<0$.

\subsection{Substeps and their control}
Each grid interval is crossed in $n$ equal substeps. The accuracy is controlled interval by interval:
\begin{enumerate}
\item Integrate the interval with $n=n_{\rm sub}$ substeps (default 4) and with $2n$ substeps, from the same starting value.
\item If the two end values differ by more than \code{rk\_tol} ($10^{-6}$ in $u$, i.e.\ a relative pressure difference of
$10^{-6}$), double $n$ and compare again ($2n$ against $4n$, \dots), up to \code{n\_sub\_max} = 64 substeps.
\item Keep the result with the larger number of substeps.
\end{enumerate}
Because the error falls by a factor 16 per doubling, the difference between two successive results is a good estimate of
the error of the coarser one, and the kept, finer result is more accurate than the tolerance. Steep temperature changes,
such as the transition region, automatically receive more substeps. Two flags are stored per column: whether any interval
needed refinement, and whether any interval still missed the tolerance at 64 substeps.

\subsection{Order of the integration for the three anchors}
\begin{itemize}
\item \code{bottom}: $u_{N_z}=\ln P_{\rm ref}$; then intervals $N_z\!\to\!N_z\!-\!1\to\dots\to1$ (upward).
\item \code{top}: $u_1=\ln P_{\rm ref}$; then $1\to2\to\dots\to N_z$ (downward).
\item \code{height}: find the interval $[z_{k_r+1},z_{k_r}]$ that contains $z_{\rm ref}$; integrate from $z_{\rm ref}$ up to
$z_{k_r}$ and on upward, and from $z_{\rm ref}$ down to $z_{k_r+1}$ and on downward.
\end{itemize}
At the anchor level the stored pressure is exactly $P_{\rm ref}$. If a substep reaches a $(T,P)$ outside the table, the
column is marked as failed and its outputs are left undefined (NaN).

\subsection{Column mass}\label{sec:mass}
Integrating Eq.~\eqref{eq:hse} from a height $z$ to the top of the grid gives
\begin{equation}
P(z)=P(z_1)+g\int_z^{z_1}\rho\,dz'\equiv P(z_1)+g\,m(z):\label{eq:colmass}
\end{equation}
\emph{the pressure at a point equals the weight of the gas above it per unit area}, plus whatever presses on the top of the
grid. The code integrates $dm/dz=-\rho$ with the same Runge--Kutta substeps as $u$ and stores $m(z_k)$ counted from the top
plane. It also stores an independent trapezoidal sum $\sum\tfrac12(\rho_k+\rho_{k+1})(z_k-z_{k+1})$ on the grid levels. The
two agree closely, and $g\,m$ should match $P-P(z_1)$; both serve as checks of the integration.

\subsection{Derived quantities}
At every grid point, with $X_k$ and $u_k$ known, the table gives $\ln\rho$, $\ln n_e$ and $\ln\chi_{500}$. From them
$P_e=n_ek_BT$, $\bar\mu=\rho k_BT/(P m_{\rm u})$ and $n_{\rm H}=\rho/(\text{mass per hydrogen nucleus})$, the latter being
fixed by the composition.

\section{Optical depth}\label{sec:tau}
\subsection{Definition}
Light of intensity $I$ crossing a path $ds$ loses $dI=-\chi I\,ds$, so $I=I_0e^{-\tau}$ with $\tau=\int\chi\,ds$, the number
of absorption lengths crossed. For light leaving vertically the optical depth of height $z$ is
\begin{equation}
\tau_{500}(z)=\tau_{\rm top}+\int_z^{z_1}\chi_{500}(z')\,dz' ,\label{eq:tau}
\end{equation}
which grows downward. For a viewing angle $\theta$ the path through each layer is longer by $1/\mu$, $\mu=\cos\theta$, and
the stored value is $\tau/\mu$ (default $\mu=1$). Light emitted at optical depth $\tau$ escapes with probability $e^{-\tau}$, so
most of the continuum light we see comes from around $\tau_{500}\approx1$, the ``visible surface''.

\subsection{Discretisation}
The integral is evaluated with the trapezoidal rule on the grid levels, starting from the top:
\begin{equation}
\tau_1=\tau_{\rm top},\qquad \tau_{k+1}=\tau_k+\tfrac12\,(\chi_k+\chi_{k+1})\,(z_k-z_{k+1}),\quad k=1\dots N_z-1 .
\end{equation}

\subsection{The gas above the grid}
The atmosphere continues above $z_1$. If $\chi_{500}$ keeps falling exponentially there,
$\chi(z)=\chi_1e^{-(z-z_1)/H_\chi}$, the optical depth of everything above the grid is
\begin{equation}
\tau_{\rm top}=\int_{z_1}^{\infty}\chi_1\,e^{-(z-z_1)/H_\chi}\,dz=\chi_1H_\chi ,\qquad
H_\chi=\frac{z_1-z_2}{\ln(\chi_2/\chi_1)} ,\label{eq:tautop}
\end{equation}
with $H_\chi$ measured from the top two levels. This is used (option \code{exp\_tail}, the default) when it makes sense:
$\chi_2>\chi_1>0$ (the extinction decreases upward) and $0<H_\chi<5000$\,km. Otherwise half a grid cell is used,
$\tau_{\rm top}=\tfrac12\chi_1(z_1-z_2)$. The options \code{half\_cell} and \code{zero} force the alternatives.

\subsection{\texorpdfstring{The height of $\tau_{500}=1$}{The height of tau500 = 1}}
In every column the code finds the two levels with $\tau_k\le1\le\tau_{k+1}$ and interpolates linearly in $\log_{10}\tau$:
\begin{equation}
z_{\tau=1}=z_k+(z_{k+1}-z_k)\,\frac{0-\log_{10}\tau_k}{\log_{10}\tau_{k+1}-\log_{10}\tau_k}.
\end{equation}

\section{Iteration and acceptance criteria}
There is no global iteration in this mode: every column is integrated once, and the result is the solution of the
discretised initial-value problem. The only accuracy criteria are local:
\begin{itemize}
\item per grid interval, $|u^{(2n)}-u^{(n)}|\le10^{-6}$ (the Runge--Kutta tolerance, Sect.~\ref{sec:num});
\item every evaluated $(T,P)$ lies inside the table envelope.
\end{itemize}
The run reports how many columns needed refinement, how many missed the tolerance, the largest number of substeps used,
how many columns used each $\tau_{\rm top}$ closure, and how many columns left the table. If any column left the table the
run ends with the status \code{eos\_failure}; otherwise it is \code{converged}. Because no reductions over columns enter the
values, the result is identical for every number of processes and threads.

\section{Three dimensions: what the result does not satisfy}\label{sec:3d}
\subsection{The residual force}
After the integration every column obeys the $z$ component of Eq.~\eqref{eq:F}. The full three-dimensional imbalance is
evaluated on the grid as a diagnostic,
\begin{equation}
F_x=-\frac{\partial P}{\partial x},\qquad F_y=-\frac{\partial P}{\partial y},\qquad F_z=-\frac{\partial P}{\partial z}-\rho g ,
\end{equation}
with three-point, second-order finite differences ($x$ and $y$ periodic, $z$ on the non-uniform grid). The layer medians
and 99th percentiles of $|F_h|/(\rho g)$ (with $F_h=\sqrt{F_x^2+F_y^2}$) and of $|F_z|/(\rho g)$ are reported: the
unbalanced force in units of the weight of the gas. $F_z$ is small (it measures only the difference between the
Runge--Kutta solution and a three-point derivative on the grid); $F_h$ is not.

\subsection{Why horizontal balance is impossible without help}
Consider two neighbouring columns at the same height, one hot and one cool. In the hot column the scale height is larger,
so above the common anchor plane the pressure falls more slowly: at the same height the hot column has a higher pressure.
There is a horizontal pressure gradient that nothing balances. Can \emph{any} pressure field make a non-magnetic
atmosphere static? Take the curl of $\F=0$. The curl of a gradient is zero, and
$\nabla\times(\rho g\zh)=g\,(\partial_y\rho,\,-\partial_x\rho,\,0)$, so
\begin{equation}
\frac{\partial\rho}{\partial x}=\frac{\partial\rho}{\partial y}=0 .
\end{equation}
A static, non-magnetic atmosphere must have a uniform density on every horizontal plane, and then also a uniform pressure.
A structured atmosphere (hot granules beside cool intergranular lanes) cannot be static: in the Sun the imbalance drives
flows, and in a static model only a magnetic force could take its place. The magnetohydrostatic modes of HSEMHS address
this; the present mode deliberately does not, and labels its output \emph{vertical HSE only}.

\section{Limitations}
\begin{itemize}
\item The horizontal force balance is not satisfied (Sect.~\ref{sec:3d}); the result is not a three-dimensional equilibrium.
\item The absolute pressure level rests entirely on $P_{\rm ref}$ and on the assumption that the anchor plane is isobaric.
A different $P_{\rm ref}$ shifts $\ln P$ of a column by nearly the same amount at all heights (exactly so for constant
$\bar\mu$), and moves the $\tau$ scale accordingly.
\item Dynamic pressure, magnetic pressure and tension are absent, so in strongly magnetised or fast-moving regions the
pressure can differ substantially from that of the original simulation.
\item The optical depth uses continuum extinction only and a model of the gas above the grid.
\end{itemize}

\section{Summary of the algorithm}
\begin{enumerate}
\item Read $T$, $x$, $y$, $z$; set $g$ and $P_{\rm ref}$; load the equation-of-state table.
\item For every column independently (in parallel): set $u=\ln P_{\rm ref}$ at the anchor; integrate $du/dz=-g\rho/P$ and
$dm/dz=-\rho$ interval by interval with RK4, $\ln T$ linear in $z$, substeps doubled until two successive results agree
to $10^{-6}$ (at most 64).
\item At every level evaluate $\rho$, $n_e$, $\chi_{500}$ from the table; derive $P_e$, $\bar\mu$, $n_{\rm H}$ and the
trapezoidal column mass.
\item Integrate $\tau_{500}$ from the top with the trapezoidal rule, starting from $\tau_{\rm top}=\chi_1H_\chi$; find
$z(\tau_{500}=1)$.
\item Evaluate the three-dimensional residual force as a diagnostic; write all fields and the per-layer statistics.
\end{enumerate}
\end{document}
